\cotefile{75}
\scanpage{26}
c) Orientation of a finite set $I$\inserted{(cardinality $n$).} The datum of $I$ is equivalent to that of a torsor
\[
 P=\operatorname{Isom}(\deleted{\unsure{[1,n]_{\mathbb Z}}}I_n,I)=\operatorname{numérot}(I)
\]
$(I_n\mathrel{\unsure{:|=}}\{m\in\mathbb N^{*}\mid m\leq n\}=\{\deleted{\unsure{0}}1,2,\ldots,n\})$ under the group
\[
 \deleted{\unsure{\mathfrak S|S}}\mathfrak S_n=\mathfrak S_{I_n}=\operatorname{Aut}(I_n).
\]

We have a canonical homomorphism ``sign''
\diagram{\begin{tikzcd}[ampersand replacement=\&,column sep=large]
\mathfrak S_n \arrow[r,"s=\mathrm{sgn}"] \& \{\pm1\} \& \text{kernel }\mathfrak A_n
\end{tikzcd}}
and by definition the set of the 2 orientations of $I$ is
\[
 \omega_I=\deleted{\unsure{\operatorname{Isom}}(}P\times_{\mathfrak S_n}\{\pm1\}\simeq P/\mathfrak A_n .
\]

\uline{Compatibilities with} a) \uline{and} b)

\uline{With} b) We have on $E_{(I,k)}=k^{I}$ a canonical scalar product
\[
 (x,y)=\sum_{i\in I}x_i y_i
\]
for which $(e_i)$ is an orthonormal basis. Thus $\deleted{E}I$ \deleted{\unsure{gives}}\inserted{\unsure{gives}} rise to a quadratic module $E(I,\unsure{x|k|X})$. This being so, we \inserted{define} a canonical bijection
\[
 \omega_{E(I,x)}\simeq\omega_I
\]
\margin{If $i_1,\mathinner{\cdotp\cdotp}i_n$ is a numbering of $I$, $e_{i_1}\mathinner{\cdotp\cdotp}e_{i_n}$ is an \inserted{\unsure{ordered}} orthonormal basis of $E$, \ldots}
by transitivity in
\diagram{\begin{tikzcd}[ampersand replacement=\&,column sep=large]
 \mathfrak S_n \arrow[r,hook,"\mathrm{incl.}"] \arrow[rr,bend right=24,"\mathrm{sgn}"] \& \mathrm{O}(\unsure{n|N},k) \arrow[r,"\det"] \& \{\pm1\}
\end{tikzcd}}

In particular, we find (compare with a).),
\[
 \omega_{E(I,\mathbb R)\tikz[baseline=0pt]{\draw[->] (0,0) .. controls (.17,-.10) and (.31,-.22) .. (.49,-.22);}}\simeq\omega_I
\]
% The arrow starts at the closing parenthesis of E(I,R) and points down/right to the two-line note at the foot.
% The lower strokes of the second line meet the bottom edge of the scan.
\margin{considered\\as \unsure{vect.|manifold} \unsure{real|via|life} \unsure{is thought of|its points}}
% archive sheet number: 25
